% ============================================================
%  05.tex --- Computer-algebra foundations: Symbolica and Spenso
% ============================================================

% ==============================================================
\chapter{Computer-Algebra Foundations: Symbolica and Spenso}
\label{ch:cas}
% ==============================================================
Why is it now possible to implement Feynman rules in Python? Previously, Python did not have a symbolic algebra engine, so it could not automatically compute Feynman rules. More importantly, it could not efficiently handle the large tensor expressions needed for quantum theory. Now, Symbolica makes this possible by providing high-performance symbolic manipulation while keeping the flexibility of Python.

Before explaining how the toolkit turns a declared model into Feynman rules in Chapter ~\ref{ch:internals}, it helps to introduce the main libraries used in the project. The toolkit represents every symbolic object, like couplings, momenta, gamma matrices, or entire vertices, with Symbolica. As mentioned in the previous chapter, indices are also important, and the Spenso library gives them meaning. Symbolica handles the expressions, Spenso explains the indices, and a small extra layer called idenso adds the physical identities for these typed indices. This chapter introduces each library, explains its role in the project, and discusses why it was chosen.




% --------------------------------------------------------------
\section{Symbolica: a modern computer-algebra core}
\label{sec:cas-symbolica}
% --------------------------------------------------------------

\textsc{Symbolica}~\cite{symbolica} is a high-performance computer algebra system (CAS) developed by Ben Ruijl. Its core is written in Rust, and it can be used through interfaces for Python, Rust, and C++. It is used in high-energy-physics research, including at CERN. It arose from the practical needs of collider physics, where symbolic computations can generate expressions containing billions of terms and requiring terabytes of memory.  Established tools, such as the \textsc{Form} family and general-purpose systems like \textsc{Mathematica}, can be difficult to maintain or are not designed for computations at this scale~\cite{symbolica,ruijl_symbolica_talk}.

Three properties of \textsc{Symbolica} are particularly relevant to this thesis and have influenced the design of the toolkit. They are introduced in the following sections.

\subsection*{Expressions, symbols, and pattern matching}

Expressions are the central objects in \textsc{Symbolica}: they are parsed, constructed, transformed, matched, differentiated, evaluated, and printed throughout the computation. A \py{Expression} is an ordinary in-process object that can be stored directly in Python data structures, compared, hashed, and passed through the toolkit without serialisation. Expressions are built from exact or approximate numbers, symbols, functions, sums, products, and powers, and are automatically normalised when created. Symbols can be created with \py{S("name")}, while complete expressions can be parsed with \py{E("...")}:



\begin{minted}{python}
from symbolica import S, E, Expression
x, y, f = S("x", "y", "f")
e = f((1 + x)**2)          # a function application, f of (1+x)^2
\end{minted}

One of the most used capability for this project is \emph{pattern matching with wildcards}. Any symbol whose name ends in an underscore is treated as a wildcard in \py{match}/\py{replace} operations, so algebraic rewriting reads almost like a regular expression for mathematics:

\begin{minted}{python}
x_, n_ = S("x_", "n_")
e = f((1 + x)**2).replace(f(x_**n_), f(x_)**n_)
print(e)   # ->  f(1 + x)**2
\end{minted}

The toolkit leans on this pervasively. As we will see in Chapter~\ref{ch:internals}, the central step of the contraction engine ( attaching an external leg's index label to an open index of a coupling) is literally one call to
\py{coupling.replace(...)}, and the user-facing helper \py{pattern_coefficient} (Section~\ref{sec:feynpy-expr-manip}) is a thin wrapper around the same wildcard machinery. Standard CAS operations are available on the same objects: \py{.expand()}, \py{.collect()}, \py{.derivative()} (with the chain and product rules applied automatically), \py{.factor()}, \py{.coefficient()}, series expansion, and even \py{solve_linear_system}.

\subsection*{Polynomial algebra: the source of the speed}

What makes \textsc{Symbolica} so useful is the highly efficient algorithms it offers for multivariate polynomial operations. These include the greatest common divisor, factorisation, and interpolation over the rationals~\cite{symbolica}. This polynomial engine is one of the main sources of its performance. For example, in Symbolica's own multivariate-GCD benchmark, a calculation that takes about $89\,\mathrm{s}$ in \textsc{Mathematica} and more than an hour in \textsc{SymPy} is completed in about $4\,\mathrm{s}$~\cite{symbolica,ruijl_symbolica_talk}.

This is important for the toolkit developed in this thesis, because the contraction engine described in Chapter~\ref{ch:internals} generates many terms by summing over permutations of the external legs. It then uses Symbolica's \py{.expand()} and \py{.collect()} methods to simplify the resulting sums and rewrite them in a compact polynomial form. Since this normalisation is performed repeatedly on potentially large expressions, the efficiency of Symbolica's polynomial engine has a direct effect on the overall performance.

\subsection*{Tensor and graph canonicalisation}

Symbolica is not limited to scalar algebra. For this project, one of its most useful features is the possibility to bring tensor expressions into a canonical form. Given a product of tensors with contracted indices, \py{Expression.canonize_tensors} renames dummy indices in a deterministic way and orders the tensor factors according to their symmetry properties. Internally, this is based on a graph-canonicalisation algorithm for tensor networks~\cite{ruijl_symbolica_talk}. In this way, expressions that differ only by dummy-index names, tensor symmetries, or the ordering of commuting factors are reduced to the same form.

In FeynPy this cannot be applied directly, because tensors are kept as typed Spenso objects during the calculation, while Symbolica expects its symmetry information on ordinary Symbolica heads. For this reason we introduced the wrapper \py{canonize_spenso_tensors}. It temporarily replaces the Spenso tensor heads by Symbolica ones, attaches the needed symmetry information, performs the canonicalisation, and then restores the original objects.

The wrapper also keeps Lorentz, spinor, colour and weak dummy indices separated. On top of this, \py{canonize_full} applies the additional identities needed in the physics pipeline. These include metric contractions, the ordering of nested flat derivatives, the antisymmetry of \(f^{abc}\), and local Jacobi reductions involving products of structure constants.

\begin{tcolorbox}[infobox, title={Why a fast CAS matters here}]
The toolkit expands covariant derivatives and field strengths into many smaller terms, and afterwards evaluates the possible Wick contractions. The number of terms can therefore grow quite fast. Having a CAS that can manipulate and canonicalise large symbolic expressions efficiently is not only convenient here, but really necessary.
\end{tcolorbox}


% --------------------------------------------------------------
\section{Spenso: typed tensor indices and tensor networks}
\label{sec:cas-spenso}
% --------------------------------------------------------------

We now turn to the role of \textsc{Spenso}. From Symbolica's point of view, an expression such as \py{G(mu, a)} is simply a function applied to two arguments. In fact it doesn't know that the index \py{mu} is a Lorentz index or that \py{a} is an adjoint colour index. Her is where \textsc{Spenso}~\cite{spenso} comes into play providng additional structure to the whole. It represents each tensor index as a typed slot, consisting of an abstract label together with a representation, and it can then match and contract compatible indices automatically.

\subsection*{Slots and representations}

The central concept in Spenso is the \py{Slot}, which combines a \py{Representation} with an abstract index label. A \py{Representation} identifies the type and dimension of the index and determines its duality and contraction rules. The toolkit uses the six representations listed in Table~\ref{tab:spenso-representations}, each created once and reused for all indices of the corresponding type.

\begin{table}[htbp]
\centering
\begin{tabular}{@{}llp{6.6cm}@{}}
\toprule
\textbf{Constructor} & \textbf{Self-dual?} & \textbf{Used for} \\
\midrule
\texttt{Representation.mink(4)}  & yes & Lorentz vector indices ($\mu,\nu,\dots$) \\
\texttt{Representation.bis(4)}   & yes & Dirac bispinor indices on fermion fields \\
\texttt{Representation.cof(3)}   & no  & $\mathrm{SU}(3)$ colour fundamental (quark) index \\
\texttt{Representation.coad(8)}  & yes & $\mathrm{SU}(3)$ colour adjoint (gluon / generator / $f^{abc}$) index \\
\texttt{Representation.cof(2)}   & no  & $\mathrm{SU}(2)_L$ weak-isospin fundamental (doublet) index \\
\texttt{Representation.coad(3)}  & yes & $\mathrm{SU}(2)_L$ weak-isospin adjoint (triplet) index \\
\bottomrule
\end{tabular}
\caption{The Spenso representations used by the toolkit. Each is wrapped by one of the toolkit's own \texttt{IndexType} constants of Section~\ref{sec:feynpy-indices} (\texttt{LORENTZ\_INDEX}, \texttt{SPINOR\_INDEX}, \texttt{COLOR\_FUND\_INDEX}, \texttt{COLOR\_ADJ\_INDEX}, \texttt{WEAK\_FUND\_INDEX}, \texttt{WEAK\_ADJ\_INDEX}).}
\label{tab:spenso-representations}
\end{table}

The \emph{self-dual} column reflects an important distinction between the representations introduced in Chapter~\ref{ch:theory}. Let's consider, for example, a colour fundamental index, carried by a quark $q^i$. We also have a colour antifundamental index, carried by an antiquark $\bar q_i$, and this belong to dual representations. In the same way therefore \py{cof(3)} is dualisable, and Spenso only contracts slots with the same abstract label when one belongs to the fundamental representation and the other to its dual. By contrast, Lorentz, bispinor, and adjoint representations are treated as self-dual in the library. Matching slots of these types may therefore be contracted directly. Once the toolkit has assigned the correct representation to each slot, Spenso can decide which repeated indices are compatible for contraction.

\subsection*{Building typed tensors}

Now we want to understand how to build and use tensors. The first thing we want is every index label wrapped into a Spenso \py{Slot}. So that an index label such as \py{mu} or \py{i} is not just name. A slot records both the label and the representation to which it belongs. Only after this typing step is the tensor converted back into a normal Symbolica \py{Expression}, so that it can be multiplied, expanded, matched, and simplified like any other symbolic factor.

In practice, constructing a typed tensor has three steps:
\begin{enumerate}
    \item choose the Spenso \py{Representation} of each index;
    \item turn each abstract label into a typed \py{Slot};
    \item call a tensor head on those slots and export the result with \py{.to_expression()}.
\end{enumerate}

For example, the Dirac gamma matrix is already part of Spenso's HEP tensor library. The toolkit wraps it as follows:
\begin{minted}{python}
def gamma_matrix(left_spinor, right_spinor, lorentz):
    return TensorName.gamma()(
        bispinor_index(left_spinor),
        bispinor_index(right_spinor),
        lorentz_index(lorentz),
    ).to_expression()
\end{minted}

Here \py{bispinor_index(left_spinor)} means
\py{Representation.bis(4)(left_spinor)}, while \py{lorentz_index(lorentz)}
means \py{Representation.mink(4)(lorentz)}. Thus the expression does not merely
contain the labels \py{left_spinor}, \py{right_spinor}, and \py{lorentz}; it
contains the information that the first two are Dirac spinor slots and the last
one is a Lorentz slot. After the call to \py{.to_expression()}, this becomes a
Symbolica expression, but the Spenso slot information remains encoded inside the
tensor arguments.

The same pattern is used for colour generators and structure constants:
\begin{minted}{python}
def gauge_generator(adj_index, fund_left, fund_right):
    return TensorName.t()(
        color_adj_index(adj_index),
        color_fund_index(fund_left),
        color_fund_index(fund_right),
    ).to_expression()

def structure_constant(a, b, c):
    return TensorName.f()(
        color_adj_index(a),
        color_adj_index(b),
        color_adj_index(c),
    ).to_expression()
\end{minted}

Metrics are built in the same spirit, but directly from the representation:
\begin{minted}{python}
def lorentz_metric(mu, nu):
    return LORENTZ.g(mu, nu).to_expression()
\end{minted}

There is one more important example to focus on, namley not every tensor needed by the toolkit is provided by Spenso's standard HEP library. A useful example is the weak doublet invariant \(\varepsilon_{ij}\). It is introduced by creating a new \py{TensorName}, and then calling it on weak-fundamental slots:
\begin{minted}{python}
WEAK_EPS2 = TensorName("weak_eps2")

def weak_eps2(i, j):
    return WEAK_EPS2(
        weak_fund_index(i),
        weak_fund_index(j),
    ).to_expression()
\end{minted}

This makes \(\varepsilon_{ij}\) a typed Spenso tensor just like the built-in gamma matrices or structure constants: its arguments know that they live in \(\mathrm{SU}(2)_L\) fundamental representation slots. The constructor does not by itself teach Symbolica that the tensor is antisymmetric. That information is added later by the canonicalisation layer, where \(\varepsilon_{ij}\) is temporarily replaced by a head marked \py{is_antisymmetric=True}. This is what allows expressions such as \(\varepsilon_{ij}\) and \(-\varepsilon_{ji}\) to be recognised as the same tensor structure.

This distinction is important. Typed Spenso tensors such as \(\gamma^\mu\), \(t^a\), \(f^{abc}\), metrics, and \(\varepsilon_{ij}\) carry representation information inside their arguments.

On the other hand in our code we have ordinary field \py{G(mu,a)} that carry indices. These are still plain Symbolica function applications. Their index meaning is stored separately in the toolkit through the field's \py{IndexType} metadata. The compiler and canonicalisation layers combine both sources of information: Spenso handles the typed tensor factors, while the toolkit metadata tells the pipeline how the indices on fields should be interpreted.

\subsection*{Tensor networks and automatic index matching}

One of the most important features of Spenso is \py{TensorNetwork}. A product of typed tensors is a network whose edges are the shared index labels; Spenso can \emph{parse such a network directly from a Symbolica expression}, match indices automatically, and contract the network either numerically, using stored components, or symbolically. Data can be stored densely (\py{DenseTensor}) or sparsely (\py{SparseTensor}), and a tensor with no data attached behaves as a purely symbolic object. The toolkit mostly uses tensors symbolically, but it also extends the HEP tensor library for the custom invariants that need explicit tensor-network evaluation.

\begin{tcolorbox}[notebox, title={Custom invariants}]
Spenso provides a built-in HEP tensor library, accessed through

\noindent \py{TensorLibrary.hep_lib()}, which contains the standard tensors such as
\(\gamma\) matrices, \(\sigma\) matrices, group generators, structure constants,
and metrics. The project adds its own tensor heads for invariants that are not part of this default library, such as the weak doublet \(\varepsilon_{ij}\), the colour Levi-Civita tensor, the symmetric \(d^{abc}\), and the charge-conjugation matrix. Only the invariants that are needed for explicit tensor-network evaluation are given numeric components in the extended library; their symmetry properties are handled separately by the canonicalisation layer.
\end{tcolorbox}

% --------------------------------------------------------------
\section{idenso: physics-aware simplification}
\label{sec:cas-idenso}
% --------------------------------------------------------------

The third piece is \textbf{idenso}, a companion to Spenso that supplies physics-aware simplification passes over typed tensor networks. Three are used by the toolkit:

\begin{itemize}
    \item \py{simplify_metrics} contracts chains of Kronecker deltas and metric tensors, e.g.\ $g^{\mu\nu} g_{\nu\rho} \to \delta^\mu{}_\rho$, and absorbs metrics contracted against other tensors;
    \item \py{simplify_color} applies $\mathrm{SU}(N)$ colour identities --- Fierz-type collapses of networks of generators $t^a$ and structure constants $f^{abc}$;
    \item \py{simplify_gamma} performs the Clifford (Dirac) algebra on chains of $\gamma$ matrices.
\end{itemize}

The toolkit composes the three into one pass, \py{simplify_invariants}. This run in a deliberately fixed order (metrics, then colour, then gamma, then metrics again) because each pass can expose fresh metric pairs left behind by the previous one. The advantage of idenso is that it does not perform any operations on subexpressions it does not recognise; therefore, the sequence comprising those steps can be safely called on any vertex factor. This pass, together with the tensor-canonicalisation step built on Symbolica's \py{canonize_tensors} (Section~\ref{sec:internals-postprocessing}), is what turns the raw permutation sum produced by the contraction engine into the compact vertex that \py{feynman_rule(...)} finally prints.

% --------------------------------------------------------------
\section{The bridge: \texttt{IndexType}}
\label{sec:cas-indextype}
% --------------------------------------------------------------

The toolkit ties the three libraries together with one small object, \py{IndexType}, discussed from the user's point of view in Section~\ref{sec:feynpy-indices}. Each \py{IndexType} pairs a single Spenso \py{Representation} with a plain string \emph{kind} (\py{"lorentz"}, \py{"spinor"}, \py{"color_fund"}, \dots) and a default label prefix. The \py{Representation} is what Spenso needs to decide contractibility; the \py{kind} string is a cheap key that the toolkit's own compiler uses to group and match indices without repeatedly inspecting the representation object.                                                                   

In summary, the division of responsibility is sharp and deliberate:
\begin{center}
\begin{tabular}{@{}lp{9.2cm}@{}}
\toprule
\textbf{Layer} & \textbf{Responsibility} \\
\midrule
\textsc{Symbolica} & Represent and manipulate symbolic expressions; pattern matching; polynomial normalisation; tensor/graph canonicalisation. \\
\textsc{Spenso}    & Give indices a representation and dimension; decide which indices contract; build and evaluate tensor networks. \\
idenso             & Apply physics identities (metrics, $\mathrm{SU}(N)$ colour, Dirac algebra) to typed tensor networks. \\
\py{IndexType}     & Bridge the above to the toolkit by pairing one representation with a stable ``kind'' key. \\
\bottomrule
\end{tabular}
\end{center}

With these foundations in place, the next chapter can describe the pipeline that consumes them: how a declared \py{Model} is progressively rewritten into elementary interaction terms and then contracted, through Symbolica and Spenso, into a Feynman rule.
